\documentclass[aspectratio=169]{beamer}
\usepackage[utf8]{inputenc}
\usepackage[english,russian]{babel}
\usepackage{cancel}
\usepackage{amssymb}
\usepackage{stmaryrd}
\usepackage{cmll}
\usepackage{graphicx}
\usepackage{amsthm}
\usepackage{tikz}
\usepackage{multicol}
\usetikzlibrary{patterns}
\usepackage{chronosys}
\usepackage{proof}
\usepackage{comment}
\usepackage{multirow}
\usepackage{diagbox}
\setbeamertemplate{navigation symbols}{}
%\usetheme{Warsaw}

\newtheorem{thm}{Теорема}[section]
\newtheorem{dfn}{Определение}[section]
\newtheorem{lmm}{Лемма}[section]
\newtheorem{exm}{Пример}[section]
\newtheorem{snote}{Пояснение}[section]

\newcommand{\divisible}%
{\mathrel{\lower.2ex%
\vbox{\baselineskip=0.7ex\lineskiplimit=0pt%
\kern6pt \hbox{.}\hbox{.}\hbox{.}}%
}}

\begin{document}

\newcommand\doubleplus{+\kern-1.3ex+\kern0.8ex}
\newcommand\mdoubleplus{\ensuremath{\mathbin{+\mkern-10mu+}}}


\begin{frame}{Разрешимость}
\begin{dfn}Язык --- множество строк\end{dfn}
\begin{dfn}Язык $L$ разрешим, если существует машина Тьюринга, которая для любого слова $w$ переходит в допускающее состояние, если $w \in L$,
и в отвергающее, если $w \notin L$.\end{dfn}
\end{frame}

\begin{frame}{Неразрешимость задачи останова}
\begin{dfn}Рассмотрим все возможные описания машин Тьюринга. Составим упорядоченные пары: описание машины Тьюринга и входная строка.
Из них выделим язык останавливающихся на данном входе машин Тьюринга.\end{dfn}
\begin{thm}Язык всех останавливающихся машин Тьюринга неразрешим\end{thm}
\begin{proof}От противного. Пусть $S(x,y)$ --- машина Тьюринга, определяющая, остановится ли машина $x$, примененная к строке $y$.\pause
\begin{center}W(x) = if (S(x,x)) \{ while (true); return 0; \} else \{ return 1; \}\end{center}\pause
Что вернёт $S(code(W),code(W))$?
\end{proof}
\end{frame}

\begin{frame}{Кодируем состояние}
\begin{enumerate}
\item внешний алфавит: $n$ 0-местных функциональных символов $q_1, \dots, q_n$; $q_\varepsilon$ --- символ-заполнитель.
\item список: $\varepsilon$ и $c(l,s)$; <<abc>> представим как $c(q_a,c(q_b,c(q_c,\varepsilon)))$.
\item положение головки: <<$ab\underline{p}q$>> как $(c(q_b,c(q_a,\varepsilon)), c(q_p,c(q_q,\varepsilon)))$.
\item внутренний алфавит: $k$ 0-местных функциональных символов $s_1, \dots, s_k$. Из них выделенные $s_s$ --- начальное и
$s_f$ --- допускающее состояние.
\end{enumerate}
\end{frame}

\begin{frame}{Достижимые состояния}
Предикатный символ $F_{x,y}(w_l,w_r,s)$: если у машины $x$ с начальной строкой $y$ состояние $s$ достижимо на строке $rev(w_l) @ w_r$. 
\pause
Будем накладывать условия: семейство формул $C_m$. \pause
Очевидно, начальное состояние достижимо:
$$C_0 := F_{x,y}(\varepsilon,y,s_s)$$
\end{frame}


\begin{frame}{Кодируем переходы}
\begin{enumerate}
\item Занумеруем переходы.\pause
\item Закодируем переход $m$: $$\langle k, s \rangle \Rightarrow \langle k', s', \rightarrow \rangle, \text{ в случае } q_k \ne q_\varepsilon$$
$C_m = \forall w_l.\forall w_r.F_{x,y}(w_l,c(q_k,w_r),s_s) \rightarrow F_{x,y}(c(q_{k'},w_l),w_r,s_{s'})$\\
(здесь требуется, чтобы под головкой находился непустой символ $q_k$, потому мы обязательно требуем, чтобы лента была
непуста)\pause

\item Переход посложнее:
$$\langle k, s \rangle \Rightarrow \langle k', s', \leftarrow \rangle, \text{ в случае } q_k \ne q_\varepsilon$$
$C_m = \forall w_l.\forall w_r.\forall t.F_{x,y}(c(t,w_l),c(q_k,w_r),s_s) \rightarrow F_{x,y}(w_l,c(t,c(q_{k'},w_r)),s_{s'}) \with
\forall w_l.\forall w_r.F_{x,y}(\varepsilon,c(q_k,w_r),s_s) \rightarrow F_{x,y}(\varepsilon,c(q_\varepsilon,c(q_{k'},w_r)),s_{s'})$\\
\pause
\item и т.п.
\end{enumerate}
\end{frame}

\begin{frame}{Итоговая формула}
$$C = C_0 \with C_1 \with \dots \with C_n$$
<<правильное начальное состояние и правильные переходы между состояниями>>\pause

\begin{thm}
Состояние $s$ со строкой $rev(w_l)@w_r$ достижимо тогда и только тогда, когда
$C \vdash F_{x,y}(w_l,w_r,s)$\pause
\end{thm}
\begin{proof}
$(\Leftarrow)$ Рассмотрим модель: предикат $F_{x,y}(w_l,w_r,s)$ положим истинным, если состояние достижимо. \pause
Это --- модель для $C$ (по построению $C_m$). \pause
Значит, доказуемость влечёт истинность (по корректности). \pause

$(\Rightarrow)$ Индукция по длине лога исполнения.
\end{proof}
\end{frame}

\begin{frame}{Неразрешимость исчисления предикатов: доказательство}
\begin{thm}Язык всех доказуемых формул исчисления предикатов неразрешим\end{thm}
Т.е. нет машины Тьюринга, которая бы по любой формуле $\alpha$ определяла, доказуема ли она.\pause
\begin{proof} Пусть существует машина Тьюринга, разрешающая любую формулу. \pause%Тогда покажем, что задача останова разрешима.\pause
На её основе тогда несложно построить некоторую машину Тьюринга, перестраивающую любую машину $S$ (с допускающим состоянием $s_f$ и входом $y$) 
в её ограничения $C$ и разрешающую формулу ИП $C \rightarrow \exists w_l.\exists w_r.F_{S,y}(w_l,w_r,s_f)$. 
Эта машина разрешит задачу останова.
\end{proof}
\end{frame}


\begin{frame}
\begin{center}\LARGE Арифметизация логики\end{center}
\end{frame}

\begin{frame}{История вопроса: арифметизация в работах Лейбница}
\begin{center}Правила, по которым можно с помощью чисел судить о правильности выводов, о формах и модусах категорических силлогизмов (1679 г.)\end{center}

\begin{itemize}
\item Любой термин --- пара взаимно простых чисел $+a-b$. Например, мудрый --- $+70-33$, благочестивый --- $+10-3$.

\item Общеутвердительное предложение (каждый $+a-b$ есть $+c-d$): $a \divisible c$ и $b \divisible d$.\\
Всякий мудрый есть благочестивый ($70 = 10\cdot 7$, $33 = 3 \cdot 11$).

\item Частноотрицательное предложение --- не верно общеутвердительное.

\item Общеотрицательное предложение --- когда $a,d$ или $b,c$ имеют общий делитель, отличный от 1:\\
Ни один благочестивый ($+10-3$) не есть несчастный ($+5-14$), так как $10=2 \cdot 5$ и $14 =2 \cdot 7$.
\end{itemize}
\end{frame}

\begin{frame}{Силлогизмы}
\begin{itemize}
\item Очевидно, силлогизм <<Barbara>> верен: $$\infer{e \divisible c, f \divisible d}{+a-b\text{ есть }+c-d\quad\quad+e-f\text{ есть }+a-b}$$

\item Неправильный силлогизм может быть верен иногда, скажем, <<AOO>> из фигуры 3:\\Всякий благочестивый ($+10-3$) есть счастливый
($+5-1$), некоторый благочестивый не есть богатый ($+8-11$), отсюда некоторый богатый не есть счастливый ($+8-11$ против $+5-1$).

\item Но для некоторых пар всё равно следствие данного неправильного силлогизма нарушается: $+12-5$, $+4-1$, $+8-11$
\end{itemize}
\end{frame}

\begin{frame}{И. Кант, <<Новое освещение первых принципов метафизического познания>>, 1755 г.}

(...) Я признаюсь, что вижу в этом суждении великого
философа лишь нечто подобное завещанию того отца у
Эзопа, который, лежа на смертном одре, поведал своим
детям, что на своем поле он зарыл клад, однако, прежде
чем успел указать им точное это место, внезапно 
скончался. Он побудил этим сыновей к неустанному раскапыванию и разрыхлению почвы,
пока они, хотя и обманутые в своих надеждах, не оказались тем не менее бесспорно
разбогатевшими благодаря тому, что повысили плодородие почвы. (...)\vspace{0.2cm}

(...) Я не стану отрицать,
что, после того как безусловно первые принципы уже 
найдены, можно кое-где применить знаковую 
комбинаторику, так как в этом случае представляется возможность 
использовать в качестве знаков и наиболее простые понятия,
а следовательно, и простейшие выражения;
однако там, где при помощи этих знаков должно быть выражено сложное
познание, вся проницательность ума оказывается как бы внезапно
повисшей над пропастью и наталкивается на неразрешимые трудности. (...)
\end{frame}

\begin{frame}{Соглашения о записи}
\begin{itemize}
\item Рассматриваем функции $\mathbb{N}^n_0\to\mathbb{N}_0$.
\item Обозначим вектор $\langle x_1, x_2, \dots, x_n\rangle$ как $\overrightarrow{x}$.
\end{itemize}
\end{frame}


\begin{frame}{Примитивно-рекурсивные функции}

\begin{dfn}[Примитивы Z, N, U, S]
\begin{enumerate}
\item Примитив <<Ноль>> ($Z$) \vspace{-0.3cm}
$$Z: \mathbb{N}_0\to\mathbb{N}_0,\ \ \ \ \ Z(x_1) = 0$$\vspace{-0.5cm}\pause
\item Примитив <<Инкремент>> ($N$) \vspace{-0.3cm}
$$N: \mathbb{N}_0\to\mathbb{N}_0,\ \ \ \ \ N(x_1) = x_1+1$$\vspace{-0.5cm}\pause
\item Примитив <<Проекция>> ($U$) — семейство функций; пусть $k,n \in \mathbb{N}_0, k \le n$\vspace{-0.3cm}
$$U^k_n: \mathbb{N}^n_0 \to \mathbb{N}_0,\ \ \ \ \ U^k_n(\overrightarrow{x}) = x_k$$\vspace{-0.5cm}\pause
\item Примитив <<Подстановка>> ($S$) --- семейство функций; пусть $g: \mathbb{N}^k_0 \to \mathbb{N}_0,\ \ f_1,\dots,f_k: \mathbb{N}^n_0 \to \mathbb{N}_0$
 \vspace{-0.3cm}
$$S\langle g,f_1,f_2,\dots,f_k \rangle (\overrightarrow{x}) = g(f_1(\overrightarrow{x}),\dots,f_k(\overrightarrow{x}))$$
\end{enumerate}
\end{dfn}
\end{frame}

\begin{frame}[fragile]{Примитивная рекурсия}
\begin{dfn}[примитив <<примитивная рекурсия>>, $R$]
Пусть $f: \mathbb{N}^n_0\to\mathbb{N}_0$ и $g: \mathbb{N}^{n+2}_0 \to\mathbb{N}_0$.
Тогда $R\langle f,g\rangle: \mathbb{N}^{n+1}_0\to\mathbb{N}_0$, причём

$$R\langle f,g\rangle(\overrightarrow{x},y)=
 \left\{\begin{array}{ll} 
  f(\overrightarrow{x}), &y=0\\
  g(\overrightarrow{x},y-1,R\langle f,g\rangle (\overrightarrow{x},y-1)), &y > 0
\end{array}\right.$$
\end{dfn}\pause

\vspace{-0.3cm}
\begin{snote}\vspace{-0.5cm}
\begin{verbatim}
                    res := f(x1...xn);
                    for yi = 0 to y-1 do
                        res := g(x1...xn,yi,res);
\end{verbatim}
\end{snote}\pause

\vspace{-0.3cm}
\begin{exm}\vspace{-0.5cm}
$$\begin{array}{ll}
    R\langle f,g\rangle(\overrightarrow{x},3) &= g(\overrightarrow{x},2,R\langle f,g\rangle(\overrightarrow{x},2)) \\\pause
 &=   g(\overrightarrow{x},2,g(\overrightarrow{x},1,R\langle f,g\rangle(\overrightarrow{x},1))) \\\pause
 &=   g(\overrightarrow{x},2,g(\overrightarrow{x},1,g(\overrightarrow{x},0,R\langle f,g\rangle(\overrightarrow{x},0)))) \\\pause
 &=  g(\overrightarrow{x},2,g(\overrightarrow{x},1,g(\overrightarrow{x},0,f(\overrightarrow{x}))))
\end{array}$$
\end{exm}

%$$R\langle f,g\rangle(\overrightarrow{x},y) = \texttt{foldl}\ (\lambda r.\lambda t.g(\overrightarrow{x},t,r))\ f(\overrightarrow{x})\ [0..y-1]$$

\end{frame}

\begin{frame}{Примитивно-рекурсивные функции}
\begin{dfn}Функция $f$ --- примитивно-рекурсивна, если может быть
выражена как композиция примитивов $Z$, $N$, $U$, $S$ и $R$.
\end{dfn}\pause

\begin{thm}
$f(x) = x+2$ примитивно-рекурсивна
\end{thm}\pause

\begin{proof}
$f = S \langle N,N \rangle$
\vspace{0.5cm}\pause

{\color{gray}
$N(x)=x+1$

$S\langle g,f\rangle (x) = g(f(x))$
}\pause

\vspace{0.5cm}
$f,g = N$\pause

$S\langle N,N\rangle (x) = N(N(x)) = (x+1)+1$
\end{proof}
\end{frame}

\begin{frame}{Примитивно-рекурсивные функции: $x+y$}
\begin{lmm}
$f(a,b) = a+b$ примитивно-рекурсивна
\end{lmm}
\begin{proof} $f = R \langle U^1_1, S\langle N, U^3_3\rangle \rangle$:
\vspace{0.2cm}\pause

{\color{gray}
$R\langle f,g\rangle(x,y)
=\left\{\begin{array}{ll} 
  f(x), &y=0\\
  g(x,y-1,R\langle f,g\rangle (x,y-1)), &y > 0
\end{array}\right.$ 

%$U^1_1(x) = x$

%$U^3_3(x,y,r) = r$
}\vspace{0.2cm}\pause

\begin{itemize}
\item База. $R\langle U^1_1, S\langle N, U^3_3\rangle \rangle(x,0) = U^1_1(x) = x$\pause

\item Переход. $R\langle U^1_1, S\langle N, U^3_3\rangle \rangle(x,y+1) = $\pause

$... = S\langle N, U^3_3\rangle (x,y,R\langle U^1_1(x), S\langle N, U^3_3\rangle\rangle (x,y) ) =$\pause

$... = S\langle N, U^3_3\rangle (x,y,x + y) = $\pause 

$... = N(x+y) = x+y+1$
\end{itemize}

\end{proof}
\end{frame}

\begin{frame}[fragile]{Какие функции примитивно-рекурсивные?}
\begin{enumerate}
\item Сложение, вычитание\pause
\item Умножение, деление\pause
\item Вычисление простых чисел\pause
\item Неформально: все функции, вычисляемые конечным числом вложенных циклов \verb!for!:

\begin{verbatim}
for (int i1 = 0; i1 < g1(x1...xn); i1++) {
    for (int i2 = 0; i2 < g2(x1...xn,i1); i2++) {
        ...
           for (int ik = 0; ik < gk(x1...xn,i1,i2...); ik++) {
               // выражение без циклов
           }
        ...
    }
}
\end{verbatim}
\end{enumerate}
\end{frame}

\begin{frame}[fragile]{Общерекурсивные функции}
\begin{dfn}
Функция --- общерекурсивная, если может быть построена при помощи
примитивов $Z$, $N$, $U$, $S$, $R$ и примитива минимизации:
$$M\langle f \rangle (x_1,x_2,\dots,x_n) = \min\{y: f(x_1,x_2,\dots,x_n,y) = 0\}$$
Если $f(x_1,x_2,\dots,x_n,y) > 0$ при любом $y$, результат не определён.
\end{dfn}\pause

%Это --- цикл \texttt{while}: 

Пример:

Пусть $f(x,y) = x-y^2$, тогда $\lceil\sqrt{x}\rceil = M\langle f\rangle (x)$

\begin{verbatim}
int sqrt(int x) {
    int y = 0;
    while (x-y*y > 0) y++;
    return y;
}
\end{verbatim}

\end{frame}

\begin{frame}{Выразительная сила}
\begin{dfn}
Функция Аккермана:
$$A(m,n) = \left\{\begin{array}{ll}
  n+1,&m = 0\\
  A(m-1,1),&m > 0, n = 0\\
  A(m-1,A(m,n-1)),&m > 0, n > 0
\end{array}\right.$$
\end{dfn}

%\begin{thm}
%Функция Аккермана --- общерекурсивная, но не примитивно-рекурсивная.
%\end{thm}\pause

\begin{exm}
  \begin{tabular}{cccccccc}
  \backslashbox{n}{m} & $0$ & $1$ & $2$ & $3$ & $4$ & \dots \\\hline
  $0$ & $1$ & $2$ & $3$ & $5$ & $13$   \\
  $1$ & $2$ & $3$ & $5$ & $13$ & $65533$ \\
  $2$ & $3$ & $4$ & $7$ & $29$ & $2^{65536}-3$ \\
  $t$ & $t+1$ & $t+2$ & $2t+3$ & $2^{t+3}-3$ & $\underbrace{2^{2^{2^{\dots^2}}}}_{t+3}-3$
  \end{tabular}
  \end{exm}
\end{frame}

\begin{frame}{Лемма о росте функции Аккермана}
\begin{multicols}{2}
\begin{dfn}$A^{(p)}(k,x) = \underbrace{A(k,A(k,A(k,\dots,A(k,x)))}_{p\ \text{раз}}$\end{dfn}
\begin{lmm}
\begin{enumerate}
\item $A(p,q) = A^{(q+1)}(p-1,1)$
\item $A^{(x+2)}(k,x) < A(k+2,x)$
\end{enumerate}\end{lmm}
\columnbreak

\color{gray}\footnotesize
$$A(m,n) = \left\{\begin{array}{ll}
  n+1,&m = 0\\
  A(m-1,1),&m > 0, n = 0\\
  A(m-1,A(m,n-1)),&m > 0, n > 0
\end{array}\right.$$

$A(0,n) = n+1$ \\

$A(2,n) = 2n + 3$

\end{multicols}

\vspace{-0.5cm}
\begin{proof}
\begin{enumerate}
\item $A(p,q) = A(p-1,A(p,q-1)) = \dots = A(p-1,A(p-1,\dots A(p,0)) = A^{(q)}(p-1,A(p,0)) = A^{(q+1)}(p-1,1)$

\item
%$\alpha_{m+2}(n) = \alpha_{m+1}^{n+1}(1) = \alpha_{m+1}(\alpha_{m+1}^n(1)) =
%\alpha_m^{\alpha_{m+1}^n(1)+1}(1) = \alpha_m^{\alpha_{m+2}(n-1)+1}(1)
%\ge \alpha_m^{\alpha_2(n-1)+1}(1) = \alpha_m^{2(n-1)+3+1}(1) =
% \alpha_m^{n+n+2}(1) = \alpha_m^{n+2} (\alpha_m^{n}(1))$

$A(k+2,x) = A(k+1,A(k+2,x-1)) = A^{(A(k+2,x-1)+1)}(k,1) \ge A^{(A(2,x-1)+1)}(k,1) = A^{(2(x-1)+3+1)}(k,1)
= A^{(2x+2)}(k,1) = A^{(x+2)}(k,A^{(x)}(k,1)) \ge A^{(x+2)}(k,A^{(x)}(0,1)) = A^{(x+2)}(k,x+1) > A^{(x+2)}(k,x)$

\end{enumerate}
\end{proof}
\end{frame}

\begin{frame}{Функция Аккермана не примитивно-рекурсивна}
\begin{thm}Пусть $f(\overrightarrow{x})$ --- примитивно-рекурсивная.
Тогда найдётся $k$, что $f(\overrightarrow{x}) < A(k,\max(\overrightarrow{x}))$\end{thm}

\begin{proof}
Индукция по структуре $f$.
\begin{enumerate}
\item $f = Z$, тогда $k = 0$, т.к. $A(0,x) = x+1 > Z(x) = 0$;
\item $f = N$, тогда $k = 1$, т.к. $A(1,x) = x + 2 > N(x) = x+1$;
\item $f = U^n_s$, тогда $k = 0$, т.к. $f(\overrightarrow{x}) \le \max(\overrightarrow{x}) < A(0,\max(\overrightarrow{x}))$;
\item $f = S\langle g,h_1,\dots,h_n\rangle$, тогда $k = k_g + \max(k_{h_1},\dots,k_{h_n}) + 2$;
%f(\overrightarrow{x}) < A(k_f,\max(\overrightarrow{x}))$, $h_t(\overrightarrow{x}) < A(k_{h_t},\max(\overrightarrow{x}))$.
%Пусть $k = \max(k_g,k_{h_1},\dots,k_{h_n})+2$;
 %$S\langle g,h_1,\dots,h_n\rangle < A(k_g, A(\max (\overrightarrow{k_h}), \max({\overrightarrow{x}))) = A^{(\max)}$
\item $f = R\langle g,h \rangle$, тогда $k = \max(k_g,k_h)+2$. %$R\langle g,h\rangle(x,v) = g^{[v]}(h(x)) < A(k_g,A(\dots A(k_g,A(k_h,x)))) \le A(x \ $
% < A(k_g,A(\dots A(k_g,A(k_h,max(\overrightarrow{x},y))))) 
%\le A^{(y}$
\end{enumerate}
\end{proof}
\end{frame}

\begin{frame}{Доказательство оценки для R}
\begin{lmm}
Пусть $f = R\langle g,h \rangle$. Тогда при $k = \max(k_g,k_h)+2$ выполнено $f(\overrightarrow{x},y) \le A^{(y+1)}(k-2,\max(\overrightarrow{x},y))$.
\end{lmm}


\begin{proof}

Индукция по $y$.
\begin{itemize}
\item База: $y = 0$. Тогда: $f(\overrightarrow{x},0) = g(\overrightarrow{x}) \le 
A(k_g,\max(\overrightarrow{x})) \le A^{(1)}(k-2,\max(\overrightarrow{x},0))$.

\item Переход: пусть $f(\overrightarrow{x},y) \le A^{(y+1)}(k-2,\max(\overrightarrow{x},y))$.
Тогда $f(\overrightarrow{x},y+1) = h(\overrightarrow{x},y,f(\overrightarrow{x},y))
\le A(k_h,\max(\overrightarrow{x},y,f(\overrightarrow{x},y))) 
\le A(k_h,\max(\overrightarrow{x},y,A^{(y+1)}(k-2,\max(\overrightarrow{x},y))) = A(k_h,A^{(y+1)}(k-2,\max(\overrightarrow{x},y)))
\le A^{(y+2)}(k-2,\max(\overrightarrow{x},y+1))$

%\le \alpha^{y+2}_{K-2}(\max(\overrightarrow{x},y)))\le$$
\end{itemize}
\vspace{-0.3cm}
\end{proof}

Заметим, что $A^{(y+1)}(k-2,\max(\overrightarrow{x},y)) \le
A^{(\max(\overrightarrow{x},y)+1)}(k-2,\max(\overrightarrow{x},y)) \le
 A^{(\max(\overrightarrow{x},y)+2)}(k-2,\max(\overrightarrow{x},y)) <
A(k,\max(\overrightarrow{x},y))$
\end{frame}

\begin{frame}{Тезис Чёрча}
\begin{dfn}
Тезис Чёрча для общерекурсивных функций: любая эффективно-вычислимая
функция $\mathbb{N}^k_0\to\mathbb{N}_0$ является общерекурсивной.
\end{dfn}
\end{frame}

\begin{frame}{Новые обозначения}
\begin{dfn}
Запись вида $\psi(\theta_1,\dots,\theta_n)$ означает
$\psi[x_1:=\theta_1,\dots,x_n:=\theta_n]$
\end{dfn}\pause

\begin{dfn}[Литерал числа]
$$\overline{a} = \left\{\begin{array}{ll} 0, &\mbox{если } a = 0\\
                (\overline{b})', &\mbox{если } a = b+1
\end{array}\right.$$
\end{dfn}\pause

Пример: пусть $\psi := x_1 = 0$. \pause Тогда $\psi(\overline{3})$ соответствует формуле $0''' = 0$

\end{frame}

\begin{frame}
\begin{center}{\LARGE Аксиоматика Пеано и формальная арифметика}\end{center}

\end{frame}


\begin{frame}{Формализуем дальше: числа}

{\itshape \hfill \begin{tabular}{r} <<Бог создал целые числа, всё остальное — дело рук человека.>>\\
                                 Леопольд Кронекер, 1886 г.\end{tabular}}\pause

\vspace{0.1cm}
Впрочем, целые числа тоже определить несложно, при наличии натуральных $\mathbb{N}_0$:
\vspace{-0.1cm}

 $$\mathbb{Z}: \dots -3, -2, -1, 0, 1, 2, 3, \dots$$\pause\vspace{-0.5cm}
\begin{itemize}
\item     $Z = \{\langle x, y \rangle\ |\ x,y\in \mathbb{N}_0\}$ \pause
\item Интуиция: $\langle x,y\rangle = x-y$\pause
\item $$\begin{array}{rcl}
      \langle a, b \rangle + \langle c, d \rangle & = & \langle a + c, b + d \rangle \\ \pause
      \langle a, b \rangle - \langle c, d \rangle & = & \langle a + d, b + c \rangle 
  \end{array}$$\pause
\item     Пусть $\langle a, b \rangle \equiv \langle c,d\rangle$, если $a + d = b + c$. Тогда $\mathbb{Z} = Z/_\equiv$\pause
\item      $0 = [\langle 0,0 \rangle],\; 1 = [\langle 1,0\rangle],\; -7 = [\langle 0,7 \rangle]$

\end{itemize}
\end{frame}

\begin{frame}{}
\begin{enumerate}
\item Рациональные ($\mathbb{Q}$).\pause

      $Q = \mathbb{Z} \times \mathbb{N}$ --- множество всех простых дробей.\pause

      $\langle p,q \rangle$ --- то же, что $\frac{p}{q}$ \pause

      $\langle p_1,q_1 \rangle \equiv \langle p_2, q_2 \rangle$, если $p_1q_2 = p_2q_1$\pause

\vspace{0.1cm}
      $\mathbb{Q} = Q/_\equiv$

\item Вещественные ($\mathbb{R}$). \pause $X = \{ A, B \}$, где $A,B \subseteq \mathbb{Q}$ --- дедекиндово сечение, если:\pause
\begin{enumerate}
\item $A\cup B = \mathbb{Q}$\pause
\item Если $a \in A$, $x \in \mathbb{Q}$ и $x \le a$, то $x \in A$\pause
\item Если $b \in B$, $x \in \mathbb{Q}$ и $b \le x$, то $x \in B$\pause
\item $A$ не содержит наибольшего.\pause
\end{enumerate}

       $\mathbb{R}$ --- множество всех возможных дедекиндовых сечений. \pause

$\sqrt 2 = \{\{ x\in\mathbb{Q}\ |\ x < 0 \vee x^2 < 2 \}, \{ x\in\mathbb{Q}\ |\ x > 0 \ \& \ x^2 > 2\}\}$
\end{enumerate}
\end{frame}


\begin{frame}{Натуральные числа: аксиоматика Пеано, 1889}\vspace{-0.5cm}
 $$\mathbb{N}: 1, 2, \dots \mbox{ или } \mathbb{N}_0: 0, 1, 2, \dots$$\vspace{-1cm}\pause
\begin{dfn}
  $N$ (или, более точно, $\langle N, 0, (')\rangle$) \emph{соответствует} аксиоматике Пеано, 
  если следующее определено/выполнено:\pause
  \begin{enumerate}
     \item Операция <<штрих>> $('): N \to N$, причём нет $a,b \in N$, что $a \ne b$, но $a' = b'$.\pause
           Если $x = y'$, то $x$ назовём следующим за $y$, а $y$ --- предшествующим $x$.\pause
     \item Константа $0 \in N$: нет $x \in N$, что $x' = 0$.\pause
     \item Индукция. Каково бы ни было свойство (<<предикат>>) $P: N \to V$, если:
           \begin{enumerate}
           \item $P(0)$
           \item При любом $x\in N$ из $P(x)$ следует $P(x')$
           \end{enumerate}
           то при любом $x \in N$ выполнено $P(x)$.
  \end{enumerate}
\end{dfn}\pause
Как построить? Например, в стиле алгебры Линденбаума:\pause
\begin{enumerate}
\item $N$ --- язык, порождённый грамматикой $\nu ::= \texttt{0}\ |\ \nu \texttt{<<'>>}$\pause
\item $0$ --- это $\texttt{<<0>>}$, $x'$ --- это $x \doubleplus \texttt{<<'>>}$
\end{enumerate}
\end{frame}


\begin{frame}{Примеры: что не соответствует аксиомам Пеано}
\begin{enumerate}
\item $\mathbb{Z}$, где $x' = x^2$\\\pause

Функция <<штрих>> не инъективна: $-3^2 = 3^2 = 9$\pause

\item Кольцо вычетов $\mathbb{Z}/{7\mathbb{Z}}$, где $x' = x+1$\\\pause

$6' = 0$, что нарушает свойства $0$\pause

\item $\mathbb{R^+}\cup\{0\}$, где $x' = x+1$\\\pause

Пусть $P(x)$ означает <<$x \in \mathbb{Z}$>>:\pause \begin{enumerate}
\item $P(0)$ выполнено: $0 \in \mathbb{Z}$.\pause
\item Если $P(x)$, то есть $x \in \mathbb{Z}$, то и $x+1 \in \mathbb{Z}$ --- так
что и $P(x')$ выполнено.\pause
\end{enumerate}
Однако $P(0.5)$ ложно.
\end{enumerate}
\end{frame}

\begin{frame}{Пример доказательства}
\begin{thm}0 единственен: если $t$ таков, что при любом $y$ 
выполнено $y' \ne t$, то $t = 0$.
\end{thm}\pause
\begin{proof}

\begin{itemize}
\item Определим $P(x)$ как <<либо $x = 0$, либо $x = y'$ для некоторого $y \in N$>>.\pause
\begin{enumerate}
\item $P(0)$ выполнено, так как $0 = 0$.\pause
\item Если $P(x)$ выполнено, то возьмём $x$ в качестве $y$: тогда для $P(x')$
будет выполнено $x' = y'$.\pause
\end{enumerate}
Значит, $P(x)$ для любого $x \in N$.\pause

\item Рассмотрим $P(t)$: <<либо $t = 0$, либо $t = y'$ для некоторого $y \in N$>>.
Но так как такого $y$ нет, то неизбежно $t = 0$.
\end{itemize}
\end{proof}
\end{frame}

\begin{frame}{Обозначения и определения}
\begin{dfn}
$1 = 0'$, $2 = 0''$, $3 = 0'''$, $4 = 0''''$, $5 = 0'''''$, $6 = 0''''''$,
$7 = 0'''''''$, $8 = 0''''''''$, $9 = 0'''''''''$
\end{dfn}\pause
\begin{dfn}\vspace{-0.3cm}
$$a + b = \left\{ \begin{array}{ll} a, & \mbox{если } b = 0\\
                                    (a + c)', & \mbox{если } b = c'
                  \end{array}\right.$$
\end{dfn}\pause\vspace{-0.3cm}
Например, $$2 + 2 = 0'' + 0'' = \pause (0'' + 0')' = \pause ((0'' + 0)')' = \pause ((0'')')' = 0'''' = 4$$\pause\vspace{-0.3cm}
\begin{dfn}\vspace{-0.3cm}
$$a \cdot b = \left\{ \begin{array}{ll} 0, & \mbox{если } b = 0\\
                                    a \cdot c + a, & \mbox{если } b = c'
                  \end{array}\right.$$
\end{dfn}
%\pause
%\begin{dfn}
%$$a^b = \left\{ \begin{array}{ll} 1, & \mbox{если } b = 0\\
%                                    a^c \cdot a, & \mbox{если } b = c'
%                  \end{array}\right.$$
%\end{dfn}
\end{frame}

\begin{frame}{Пример: коммутативность сложения (лемма 1)}

\begin{multicols}{2}
\begin{lmm}[1]
$a + 0 = 0 + a$
\end{lmm} \pause
{\color{gray}
$$a + b = \left\{ \begin{array}{ll} a, & \mbox{если } b = 0\\
                                    (a + c)', & \mbox{если } b = c'
                  \end{array}\right.$$} \end{multicols}\pause
\vspace{-1.5cm}\begin{proof} Пусть $P(x)$ --- это $x + 0 = 0 + x$.
\begin{enumerate}\pause
\item Покажем $P(0)$. $0 + 0 = 0 + 0$ \pause
\item Покажем, что если $P(x)$, то $P(x')$. Покажем $P(x')$, то есть $x' + 0 = \dots$ \pause

\begin{center}
  \begin{tabular}{lrl} %Равенство & обоснование \\\hline
                      $\dots = x'$ & $a=x',b=0$: & $x' + 0 \Rightarrow x'$ \\\pause
                      $\dots = (x)'$ & \\ \pause
                      $\dots = (x + 0)'$ & $a=x,b=0$: &$(x + 0) \Leftarrow (x)$ \\ \pause
                      $\dots = (0 + x)'$ & $P(x)$: &$(x + 0) \Rightarrow (0 + x)$ \\ \pause
                      $\dots = 0 + x'$ & $a=0,b=x'$: &$0 + x' \Leftarrow (0 + x)'$
   \end{tabular}
\end{center}
\end{enumerate}\pause
Значит, $P(a)$ выполнено для любого $a \in N$.
\end{proof}
\end{frame}

\begin{frame}{Пример: коммутативность сложения (завершение)}
\begin{lmm}[2]
$a + b' = a' + b$
\end{lmm}\pause
\begin{proof} $P(x)$ --- это $a + x' = a' + x$\pause
\begin{enumerate}
\item $a + 0' = (a + 0)' = (a)' = a' = a' + 0$\pause
\item Покажем, что $P(x')$ следует из $P(x)$: $a + x'' = (a + x')' = (a' + x)' = a' + x'$
\end{enumerate}
\end{proof}\pause

\begin{thm}
$a + b = b + a$
\end{thm}\pause
\begin{proof}[Доказательство индукцией по $b$: $P(x)$ --- это $a + x = x + a$]
\begin{enumerate}
\item $a + 0 = 0 + a$ (лемма 1)\pause
\item $a + x' = (a + x)' = (x + a)' = x + a' = x' + a$
\end{enumerate}
\end{proof}
\end{frame}

\begin{frame}{Уточнение исчисления предикатов}
\begin{itemize}
\item Пусть требуется доказывать утверждения про равенство. Введём $E(p,q)$ --- предикат <<равенство>>.\pause
\item Однако $\not\vdash E(p,q)\to E(q,p)$: если $D = \{0,1\}$ и $E(p,q) ::= (p>q)$,
то $\not\models E(p,q)\to E(q,p)$.\pause
\item Конечно, можем указывать $\forall p.\forall q.E(p,q)\to E(q,p) \vdash \varphi$.\pause
\item Но лучше добавим аксиому $\forall p.\forall q.E(p,q)\to E(q,p)$.\pause
\item Добавив необходимые аксиомы, получим \emph{теорию первого порядка}.
\end{itemize}
\end{frame}

\begin{frame}{Теория первого порядка}
\begin{dfn}
Теорией первого порядка назовём исчисление предикатов с дополнительными (<<нелогическими>>
или <<математическими>>):
\begin{itemize}
\item предикатными и функциональными символами;
\item аксиомами.
\end{itemize}

Сущности, взятые из исходного исчисления предикатов, назовём \emph{логическими}
\end{dfn}
\end{frame}

\begin{frame}{Порядок логики/теории}
\begin{tabular}{llll}
Порядок & Кванторы & Формализует суждения\dots & Пример\\\hline
нулевой & запрещены & об отдельных значениях & И.В.\\
первый & по предметным переменным & о множествах & И.П.\\
%    &   & $S = \{ t\ |\ \psi[x := t] \}$ \\
    &   \multicolumn{2}{l}{\color{olive}$\{2,3,5,7,\dots\} = \{ t\ |\ \forall p.\forall q.(p \ne 1 \with q \ne 1) \rightarrow (t \ne p\cdot q)\}$}\\
второй & по предикатным переменным & о множествах множеств & Типы\\
    &   \multicolumn{2}{l}{\color{olive}$S = \{ \{t\ |\ P(t)\}\ |\ \varphi[p := P] \}$}\\
 & \dots 
\end{tabular}
\pause\vspace{0.3cm}

\begin{exm}[логики 2 порядка]\upshape
\begin{tabular}{ll}
$\alpha\rightarrow\beta\rightarrow\alpha$ (сх. акс. 1) & $\forall a.\forall b.a \rightarrow b \rightarrow a$ \vspace{0.1cm}\\
\texttt{let rec map f l = match l with} & $map: \forall a.\forall b.(a \rightarrow b) \rightarrow a\texttt{ list} \rightarrow b\texttt{ list}$ \\
\texttt{| [] -> []} \\
\texttt{| l1::ls -> f l1 :: map f ls}\vspace{0.1cm}\\
\texttt{map ((+) 1) [1;2;3] = [2;3;4]}
\end{tabular}
\end{exm}

\end{frame}

\begin{frame}{Формальная арифметика}
\begin{dfn}
Формальная арифметика --- теория первого порядка, со следующими добавленными нелогическими \dots
\begin{itemize}
\item двухместными функциональными символами $(+)$, $(\cdot)$; одноместным функциональным символом $(')$, 
нульместным функциональным символом $0$;\pause
\item двухместным предикатным символом $(=)$;\pause
\item восемью нелогическими \emph{аксиомами}:\vspace{0.1cm}
\begin{tabular}{ll}
(A1) $a=b \to a=c \to b=c$             &(A5) $a+0 = a$                     \\
(A2) $a=b \to a'=b'$                   &(A6) $a+b' = (a+b)'$               \\
(A3) $a'=b' \to a=b$                   &(A7) $a\cdot 0 = 0$                \\
(A4) $\neg a' = 0$                     &(A8) $a\cdot b' = a \cdot b + a$
\end{tabular}\pause
\item нелогической схемой аксиом индукции $\psi[x:=0]\with(\forall x.\psi\to \psi[x:=x'])\to \psi$ с метапеременными $x$ и $\psi$.
\end{itemize}
\end{dfn}
\end{frame}

\begin{frame}{Докажем, что $a=a$}
\small
Пусть $\top ::= 0=0\to 0=0 \to 0=0$, тогда:\pause

\begin{tabular}{lll}
(1) & $a=b\to a=c \to b=c$ & (Акс. А1)\\
(2) & $(a=b\to a=c \to b=c) \to \top \to (a=b\to a=c \to b=c)$ & (Сх. акс. 1)\\
(3) & $\top \to (a=b\to a=c \to b=c)$ & (M.P. 1, 2)\\
(4) & $\top \to (\forall c.a = b\to a = c \to b = c)$ & (Введ. $\forall$)\\\pause
(5) & $\top \to (\forall b.\forall c.a = b\to a = c \to b = c)$ & (Введ. $\forall$)\\
(6) & $\top \to (\forall a.\forall b.\forall c.a = b\to a = c \to b = c)$ & (Введ. $\forall$)\\\pause
(7) & $\top$ & (Сх. акс 1)\\
(8) & $(\forall a.\forall b.\forall c.a = b\to a = c \to b = c)$ & (M.P. 7, 6)\\\pause
(9) & $(\forall a.\forall b.\forall c.a = b\to a = c \to b = c) \to $\\
    & $\to (\forall b.\forall c.a+0 = b\to a+0 = c \to b = c)$ & (Сх. акс. 11)\\\pause
(10) & $\forall b.\forall c.a+0 = b\to a+0 = c \to b = c$ & (M.P. 8, 9)\\\pause
(12) & $\forall c.a+0 = a\to a+0 = c \to a = c$ & (M.P. 10, 11)\\\pause
(14) & $a+0 = a\to a+0 = a \to a = a$ & (M.P. 12, 13)\\\pause
(15) & $a+0 = a$ & (Акс. А5)\\
(16) & $a+0 = a \to a = a$ & (M.P. 15, 14)\\
(17) & $a = a$ & (M.P. 15, 16)
\end{tabular}
\end{frame}

\begin{comment}

\begin{frame}{Выразимость отношений в Ф.А.}
\begin{dfn}
Будем говорить, что отношение $R\subseteq \mathbb{N}^n_0$ выразимо в ФА, 
если существует формула $\rho$, что:
\begin{enumerate}
\item если $\langle a_1,\dots,a_n \rangle \in R$, то $\vdash \rho(\overline{a_1},\dots,\overline{a_n})$
\item если $\langle a_1,\dots,a_n \rangle \notin R$, то $\vdash \neg\rho(\overline{a_1},\dots,\overline{a_n})$
\end{enumerate}
\end{dfn}\pause

\begin{thm} отношение <<равно>> выразимо в Ф.А.:
$R = \{ \langle x,x \rangle\ |\ x \in \mathbb{N}_0 \}$
\end{thm}\pause

\begin{proof}
Пусть $\rho := x_1=x_2$. \pause Тогда:
\begin{itemize}
\item $\vdash p = p$ при $p := \overline{k}$ при всех $k \in \mathbb{N}_0$: \pause $\vdash 0=0$, $\vdash 0'=0'$, $\vdash 0''=0''$, ... \pause
\item $\vdash \neg p = q$ при $p := \overline{k}$, $q := \overline{s}$ при всех $k,s \in \mathbb{N}_0$ и $k \ne s$.

\pause $\vdash \neg 0 = 0'$, $\vdash \neg 0 = 0''$, $\vdash \neg 0''' = 0'$, ...
\end{itemize}
\end{proof}

\end{frame}

\begin{frame}{Представимость функций в Ф.А.}
\begin{dfn}
Будем говорить, что функция $f: \mathbb{N}^n_0\to\mathbb{N}_0$ представима в ФА, 
если существует формула $\varphi$, что:
\begin{enumerate}
\item если $f(a_1,\dots,a_n) = u$, то $\vdash \varphi(\overline{a_1},\dots,\overline{a_n},\overline{u})$
\item если $f(a_1,\dots,a_n) \ne u$, то $\vdash \neg\varphi(\overline{a_1},\dots,\overline{a_n},\overline{u})$
\item для всех $a_i \in \mathbb{N}_0$ выполнено $\vdash (\exists x.\varphi(\overline{a_1},\dots,\overline{a_n},x)) \with 
   (\forall p.\forall q.\varphi(\overline{a_1},\dots,\overline{a_n},p)\with \varphi(\overline{a_1},\dots,\overline{a_n},q)\rightarrow p=q)$
\end{enumerate}
\end{dfn}
\end{frame}

\begin{frame}{Соответствие рекурсивных и представимых функций}
\begin{thm}Любая рекурсивная функция представима в Ф.А.\end{thm}\pause

\begin{thm}Любая представимая в Ф.А. функция рекурсивна.\end{thm}
\end{frame}

\begin{frame}{Примитивы Z, N, U представимы в Ф.А.}
\begin{thm}
Примитивы $Z$, $N$ и $U^k_n$ представимы в Ф.А.
\end{thm}\pause

\begin{proof}
\begin{itemize}
\item $\zeta(x_1,x_2) := x_2=0$, \pause формальнее: $\zeta(x_1,x_2) := x_1=x_1 \with x_2=0$\pause
\item $\nu(x_1,x_2) := x_2=x_1'$\pause
\item $\upsilon(x_1,\dots,x_n,x_{n+1}) := x_k = x_{n+1}$ \pause \vspace{0.1cm}

формальнее: 
     $\upsilon(x_1,\dots,x_n,x_{n+1}) := (\underset{i\ne k,n+1}{\with} x_i=x_i) \with x_k = x_{n+1}$
\end{itemize}
\end{proof}
\end{frame}

\begin{frame}{Примитив S представим в Ф.А.}

$$S\langle f,g_1,\dots,g_k\rangle(x_1,\dots,x_n) = f(g_1(x_1,\dots,x_n),\dots,g_k(x_1,\dots,x_n))$$\pause
\begin{thm}
Пусть функции $f,g_1,\dots,g_k$ представимы в Ф.А. Тогда $S\langle f,g_1,\dots,g_k \rangle$ представима в Ф.А.
\end{thm}\pause

\begin{proof}
Пусть $f$, $g_1$, ..., $g_k$ представляются формулами $\varphi$, $\gamma_1$, ..., $\gamma_k$. \pause

Тогда 
$S\langle f,g_1,\dots,g_k\rangle$ будет представлена формулой
$$\exists g_1.\dots.\exists g_k.\varphi(g_1,\dots,g_k,x_{n+1})\with\gamma_1(x_1,\dots,x_n,g_1)\with\dots\with\gamma_k(x_1,\dots,x_n,g_k)$$
\end{proof}

\end{frame}

\begin{frame}{$\beta$-функция Гёделя}
Задача: закодировать последовательность натуральных чисел произвольной длины.\pause
\begin{dfn}$\beta$-функция Гёделя: $\mathcal{\beta}(b,c,i) := b \% (1 + (i+1) \cdot c)$\\
Здесь (\%) --- остаток от деления.
\end{dfn}\pause
\begin{thm}$\beta$-функция Гёделя представима в Ф.А. формулой
$$\hat{\beta}(b,c,i,d) := \exists q.(b = q \cdot (1 + c \cdot (i+1)) + d) \& (d < 1 + c \cdot (i+1))$$
\end{thm}\pause\vspace{-0.3cm}
Деление $b$ на $x$ с остатком: найдутся частное $(q)$ и остаток $(d)$, что
$b = q\cdot x + d$ и $0 \le d < x$.
\pause
\begin{thm}Если $a_0, \dots, a_n \in \mathbb{N}_0$, то найдутся такие $b,c \in \mathbb{N}_0$, что
$a_i = \beta(b,c,i)$
\end{thm}
\end{frame}

\begin{frame}{Доказательство свойства $\beta$-функции}
\begin{thm}Китайская теорема об остатках (вариант формулировки): если $u_0, \dots, u_n$ --- попарно 
взаимно просты, и $0 \le a_i < u_i$, то существует такой $b$, что $a_i = b \% u_i$.
\end{thm}\pause
Положим $c = \max(a_0,\dots,a_n,n)!$ и $u_i = 1+c\cdot(i+1)$.\pause
\begin{itemize}
\item $\text{НОД}(u_i,u_j) = 1$, если $i \ne j$.\pause

Пусть $p$ --- простое, $u_i\divisible p$ и $u_j\divisible p$ ($i < j$). \pause
Заметим, что $u_j-u_i = c \cdot (j-i)$. Значит, $c\divisible p$ или $(j-i)\divisible p$. \pause
Так как $j-i \le n$, то $c \divisible (j-i)$, потому если и $(j-i)\divisible p$, всё равно $c \divisible p$. \pause
Но и $(1+c\cdot(i+1))\divisible p$, отсюда $1 \divisible p$ --- что невозможно. \pause
\item $0 \le a_i < u_i$.\pause
\end{itemize}
Условия китайской теоремы об остатках выполнены и найдётся $b$, что \vspace{-0.2cm}
$$a_i = b \% (1 + c\cdot(i+1)) = \beta(b,c,i)$$
\qed
\end{frame}

\begin{frame}{Примитив <<примитивная рекурсия>> представим в Ф.А.}

Пусть $f: \mathbb{N}^n_0\to\mathbb{N}_0$ и $g:\mathbb{N}^{n+2}_0\to\mathbb{N}_0$ представлены формулами $\varphi$ и $\gamma$. 

Зафиксируем $x_1, \dots, x_n, y \in \mathbb{N}_0$.\pause\vspace{0.3cm}

\begin{tabular}{lll}
Шаг вычисления & Об. & Утверждение в Ф.А.\\
$R\langle f,g\rangle (x_1,\dots,x_n,0) = f(x_1,\dots,x_n)$ & $a_0$ & $\vdash \varphi(\overline{x_1},\dots,\overline{x_n},\overline{a_0})$\\ \pause
$R\langle f,g\rangle (x_1,\dots,x_n,1) = g(x_1,\dots,x_n,0,a_0)$ & $a_1$ & $\vdash \gamma(\overline{x_1},\dots,\overline{x_n},0,\overline{a_0},\overline{a_1})$\\ \pause
$\dots$\\
$R\langle f,g\rangle (x_1,\dots,x_n,y) = g(x_1,\dots,x_n,y-1,a_{y-1})$ & $a_y$ & $\vdash \gamma(\overline{x_1},\dots,\overline{x_n},\overline{y-1},\overline{a_{y-1}},\overline{a_y})$
\end{tabular}\pause\vspace{0.3cm}

По свойству $\beta$-функции, найдутся $b$ и $c$, что
$\beta (b,c,i) = a_i$ для $0 \le i \le y$.\pause

\begin{thm}Примитив $R\langle f,g\rangle$ представим в Ф.А. формулой $\rho(x_1,\dots,x_n,y,a)$:\vspace{-0.1cm}
$$\begin{array}{l}\exists b. \exists c. (\exists a_0. \hat{\beta} (b,c,0,a_0) \& \varphi (x_1,...x_n, a_0)) \\
       \&\;\;\;\;\forall k.k < y \rightarrow \exists d . \exists e . \hat{\beta} (b,c,k,d) \& \hat{\beta} (b,c,k',e) \& \gamma (x_1,..x_n,k,d,e) \\
       \&\;\;\;\;\hat{\beta} (b,c,y,a) 
\end{array}$$
\end{thm}
\end{frame}

\begin{frame}{Представимость рекурсивных функций в Ф.А.}
\begin{thm}Пусть функция $f:\mathbb{N}^{n+1}_0 \to \mathbb{N}_0$ представима в Ф.А.
формулой $\varphi(x_1,\dots,x_{n},y,r)$. Тогда примитив $M\langle f\rangle$ представим в Ф.А. 
формулой $$\mu(x_1,\dots,x_n,y) := \varphi(x_1,\dots,x_n,y,0) \with \forall u.u < y \to \neg\varphi(x_1,\dots,x_n,u,0)$$
\end{thm}\pause

\begin{thm}Если $f$ --- рекурсивная функция, то она представима
в Ф.А.
\end{thm}
\begin{proof}Индукция по структуре $f$.
\end{proof}
\end{frame}

\begin{frame}{Рекурсивность представимых функций в Ф.А.}
Фиксируем $f$ и $x_1, x_2, \dots, x_n$. Обозначим $y = f(x_1,x_2,\dots,x_n)$. \pause
По представимости нам известна $\varphi$, что $\vdash \varphi(\overline{x_1},\overline{x_2},\dots,\overline{x_n},\overline{y})$. \pause
Давайте просто переберём все результаты и доказательства!\pause

\begin{enumerate}
\item Закодируем доказательства натуральными числами.\pause
\item Напишем рекурсивную функцию, проверяющую доказательства на корректность.\pause
\item Параллельный перебор значений и доказательств: $s = 2^y \cdot 3^p$. Переберём все $s$, по $s$ получим $y$ и $p$.
Проверим, что $p$ --- код доказательства $\vdash \varphi(\overline{x_1},\overline{x_2},\dots,\overline{x_n},\overline{y})$.
\end{enumerate}
\end{frame}

\begin{frame}{Гёделева нумерация}

\begin{enumerate}
\item Отдельный символ.\vspace{0.2cm}

\begin{tabular}{lc|lc||cll}
Номер & Символ & Номер & Символ & Имя & $k,n$ & Гёделев номер\\
3 & ( &               17 & $\&$ &  0 & $0,0$ & $27 + 6$\\
5 & ) &               19 & $\forall$ & $(')$ & $0,1$ & $27 + 6 \cdot 3$\\
7 & , &               21 & $\exists$ & $(+)$ & $0,2$ & $27 + 6 \cdot 9$\\
9 & . &               23 & $\vdash$ & $(\cdot)$ & $1,2$ & $27 + 6 \cdot 2 \cdot 9$\\
11 & $\neg$ &         $25 + 6\cdot k$ & $x_k$ & $(=)$ & $0,2$ & $29 + 6 \cdot 9$\\
13 & $\rightarrow$ &  $27 + 6\cdot 2^k \cdot 3^n$ & $f_k^n$\\
15 & $\vee$ &         $29 + 6\cdot 2^k \cdot 3^n$ & $P_k^n$ 
\end{tabular}\pause

\item Формула. $\phi \equiv s_0s_1\dots s_{n-1}$. Гёделев номер: $\ulcorner\phi\urcorner = 2^{\ulcorner s_0\urcorner}\cdot 3^{\ulcorner s_1\urcorner} 
\cdot \dots \cdot p_{n-1}^{\ulcorner s_{n-1}\urcorner}$.\pause

\item Доказательство. $\Pi = \delta_0\delta_1\dots\delta_{k-1}$, его гёделев номер: $\ulcorner\Pi\urcorner =
2^{\ulcorner \delta_0\urcorner}\cdot 3^{\ulcorner \delta_1\urcorner} \cdot \dots \cdot p_{k-1}^{\ulcorner \delta_{k-1}\urcorner}$
\end{enumerate}
\end{frame}

\begin{frame}{Проверка доказательства на корректность}
\begin{thm}Следующая функция рекурсивна:
$$\text{proof}(f,x_1,x_2,\dots,x_n,y,p) = \left\{\begin{array}{ll} 
0, & \mbox{если} \vdash\phi(\overline{x_1},\overline{x_2},\dots,\overline{x_n},\overline{y}),\\
   & p \mbox{ --- гёделев номер вывода}, f = \ulcorner\phi\urcorner \\ 
1, & \mbox{иначе}
\end{array}\right.$$\vspace{-0.1cm}\pause
\end{thm}
\begin{proof}[Идея доказательства]
\begin{enumerate}
%\item По представимости, $\vdash\phi(\overline{x_1},\overline{x_2},\dots,\overline{x_n},\overline{y})$ 
%эквивалентно $f(x_1,x_2,\dots,x_n)=y$.\pause
\item Проверка доказательства вычислима.\pause
\item Согласно тезису Чёрча, любая вычислимая функция
вычислима с помощью рекурсивных функций.
\end{enumerate}
\end{proof}
\end{frame}

\begin{frame}{Перебор доказательств}
\begin{lmm}Следующие функции рекурсивны:
\begin{enumerate}
\item Функции $\text{plog}_k(n) = \max\{p: n \divisible k^p\}$, $\text{fst}(x) = \text{plog}_2(x)$ и $\text{snd}(x) = \text{plog}_3(x)$.\pause
\item Числовые литералы: $\overline{k}: \mathbb{N}_0 \rightarrow \mathbb{N}_0$, $\overline{k}(x) = k$.\pause
\end{enumerate}
\end{lmm}
\begin{thm}Если $f: \mathbb{N}^n_0\rightarrow\mathbb{N}_0$, и $f$ представима в Ф.А. формулой $\varphi$, то $f$ --- рекурсивна.
\end{thm}\pause
\begin{proof}
Пусть заданы $x_1,x_2,\dots,x_n$. Ищем $\langle y, p\rangle$, 
что $\text{proof}(\ulcorner\varphi\urcorner,x_1,x_2,\dots,x_n,y,p)=1$, \pause
напомним: $y = f(x_1,x_2,\dots,x_n)$, $p = \ulcorner\Pi\urcorner$,
$\Pi$ --- доказательство $\varphi(\overline{x_1},\overline{x_2},\dots,\overline{x_n},\overline{y})$. \pause
$$f = S \langle \text{fst}, M\langle S \langle \text{proof}, \overline{\ulcorner\varphi\urcorner}, U^1_{n+1}, U^2_{n+1}, \dots, U^n_{n+1}, 
  S \langle \text{fst},U^{n+1}_{n+1}\rangle, S\langle \text{snd}, U^{n+1}_{n+1}\rangle \rangle \rangle \rangle$$
\end{proof}
\end{frame}

\end{comment}

\end{document}
